\documentclass[12pt]{article}
\usepackage[utf8]{inputenc}
\usepackage[T1]{fontenc}
\usepackage[spanish]{babel}
\usepackage[letterpaper, margin=2.5cm]{geometry}
\usepackage{longtable}
\usepackage{array}

\begin{document}

\section*{Tabla 1: Artículos de Actividad en Clase}

\begin{longtable}{|>{\raggedright\arraybackslash}p{3cm}|>{\raggedright\arraybackslash}p{4.2cm}|>{\raggedright\arraybackslash}p{4.2cm}|>{\raggedright\arraybackslash}p{4.2cm}|}
\hline
\textbf{Criterio} & \textbf{Artículo 1} & \textbf{Artículo 2} & \textbf{Artículo 3} \\
\hline
\endfirsthead

\hline
\textbf{Criterio} & \textbf{Artículo 1} & \textbf{Artículo 2} & \textbf{Artículo 3} \\
\hline
\endhead

\textbf{Título} & Cambio climático y responsabilidad individual: Una respuesta a Johnson & Obligaciones éticas en una tragedia de los bienes comunes & Clima, acción colectiva y obligaciones morales individuales \\
\hline
\textbf{Autores y Afiliación} & Marion Hourdequin, es profesora de Filosofía en Colorado College, con trayectoria e investigación en ética ambiental. & Baylor L. Johnson, profesor de Filosofía en St. Lawrence University. & Marion Hourdequin, profesora de Filosofía en Colorado College. \\
\hline
\textbf{Pertinencia} & Considero que este texto es muy importante porque se le planta de frente a esa visión pesimista de que las acciones que uno hace a diario en la casa no sirven de nada frente a un problema tan gigante como el cambio climático. & Este artículo es interesante porque evita que toda la culpa del problema caiga encima del consumidor, poniendo el foco donde debe estar, en la responsabilidad de las instituciones y del Estado. & Este escrito me parece relevante para pensar en la coherencia de nuestras acciones, porque no podemos andar pidiéndole cambios drásticos al gobierno o a las empresas si nosotros en la casa no asumimos ni un poco de responsabilidad. \\
\hline
\textbf{Introducción} & El escrito le responde directamente a Johnson para defender que lo que hacemos individualmente sí tiene peso, buscando replantear la famosa "tragedia de los bienes comunes", al final demuestra que cuando uno reduce su huella personal, no solo está haciendo un cambio chiquito, sino les manda un mensaje ético y claro a los demás y los motiva a actuar también. & Johnson busca responder si de verdad estamos obligados a hacer sacrificios personales cuando el problema es colectivo y gigante, su planteamiento concluye que nuestra verdadera responsabilidad no es sufrir solos, sino presionar y apoyar acuerdos que apliquen para todo el mundo por igual. & Argumenta que sí tenemos el deber moral de bajar nuestras emisiones, no por salvar el planeta solos, sino por un tema de integridad personal y de respeto con las relaciones que construimos con los demás, concluye que, para tener la moral de exigir normas públicas primero debemos ser coherentes en la vida privada de cada uno. \\
\hline
\textbf{Importancia y aporte} & Me parece clave porque nos muestra que el compromiso individual y las políticas públicas tienen que ir cogidos de la mano, no es que uno reemplace al otro ni que sean cosas opuestas. & Aporta una mirada bastante realista, porque nos recuerda que no logramos mucho cambiando hábitos a puerta cerrada si al mismo tiempo no estamos exigiéndole a leyes e instituciones firmes. & Es un aporte importante porque nos recuerda que la protesta y la exigencia política cogen más fuerza y credibilidad cuando están respaldadas con el propio ejemplo. \\
\hline
\textbf{Métodos: Datos y análisis} & Analizó conductas de todos los días, como usar la bicicleta o ponerse a reciclar, y cómo eso impacta en el círculo cercano. & Usó escenarios teóricos para analizar cómo se comporta un grupo de gente cuando comparte un recurso común pero no hay reglas claras de por medio. & Analizó conceptualmente cómo nos conectamos con la comunidad y de qué manera lo que hacemos a diario termina tocando esos vínculos sociales. \\
\hline
\textbf{Métodos: Justificación} & Prefirió irse por este camino porque la teoría de juegos tradicional siempre asume que la gente actúa pura y exclusivamente por egoísmo, mientras que los enfoques tipo Ostrom reconocen que a las personas también nos mueven los valores y las normas sociales. & Uso este modelo porque argumenta que, si no existe una ley general, el esfuerzo que uno hace con tanto cariño termina siendo aplastado por el consumo desmedido de los demás. & Eligió esa forma de trabajo porque sostiene que los humanos no somos seres aislados, sino seres interconectados que nos influenciamos mutuamente todo el tiempo. \\
\hline
\textbf{Resultados} & Muestra que lo que hacemos en el día a día no es del todo invisible, sino que se vuelve una especie de mensaje o señal que termina influyendo en la forma de actuar de la gente que nos rodea. & Demuestra a nivel teórico que hacer sacrificios aislados sin obligar a los demás a cumplir las mismas reglas no soluciona el problema de fondo. & Deja claro que cuando uno dice una cosa, pero hace otra en la vida diaria, pierde totalmente la autoridad moral a la hora de exigirle cambios reales a las instituciones. \\
\hline
\textbf{Tablas, leyendas y figuras} & Al ser un trabajo de reflexión filosófica pura y dura, no trae gráficas ni tablas estadísticas. & Tampoco incluye tablas ni esquemas gráficos, ya que su desarrollo es de análisis teórico. & Es un texto analítico y de reflexión moral, por lo que no lleva tablas ni esquemas gráficos. \\
\hline
\textbf{Discusión: Aporte al tema} & Nos ayuda a entender que las leyes y las grandes políticas no se sostienen solas en el aire, sino que necesitan de la confianza y los valores de la propia comunidad para que realmente funcionen. & Nos invita a cambiar el pensamiento y pasar del simple consumo responsable, al activismo político y ciudadano. & Conecta lo que hacemos a puerta cerrada en nuestra intimidad con nuestro papel como ciudadanos activos en la sociedad. \\
\hline
\textbf{Próxima investigación} & Sería interesante investigar cuánta gente en un barrio o comunidad necesita adoptar hábitos sostenibles antes de que las autoridades locales hagan el deber de aprobar una normativa ambiental obligatoria. & Valdría la pena investigar qué clase de incentivos o argumentos harían que los ciudadanos voten a favor de leyes ambientales mucho más estrictas, aunque impliquen cambios difíciles en su vida diaria. & Se podría analizar cómo la percepción de coherencia en los líderes ambientales influye en la disposición de la gente a seguir sus propuestas o apoyarlos en las urnas. \\
\hline
\textbf{Conclusión} & Actuar bien por decisión propia no va a solucionar el problema global por sí solo, pero sí ayuda a construir ese piso de confianza necesario para que las leyes e instituciones puedan dar resultados. & Frente a un desastre global, nuestra principal obligación ética no es andar haciendo sacrificios solos, sino impulsar y exigir reglas colectivas que nos obliguen a todos a cumplir por igual dejando la desigualdad de lado. & Por integridad nos toca reducir la huella personal, no porque con eso arreglemos el clima por arte de magia, sino porque le da un respaldo ético indestructible a cualquier cambio político que salgamos a exigir. \\
\hline
\end{longtable}

\newpage

\section*{Tabla 2: Matriz de Lectura - Una biología sin individuos ontológicos}

\begin{longtable}{|>{\raggedright\arraybackslash}p{3cm}|>{\raggedright\arraybackslash}p{4.2cm}|>{\raggedright\arraybackslash}p{4.2cm}|>{\raggedright\arraybackslash}p{4.2cm}|}
\hline
\textbf{Criterio} & \textbf{Artículo 1} & \textbf{Artículo 2} & \textbf{Artículo 3} \\
\hline
\endfirsthead

\hline
\textbf{Criterio} & \textbf{Artículo 1} & \textbf{Artículo 2} & \textbf{Artículo 3} \\
\hline
\endhead

\textbf{Título} & Los límites del yo: inmunología e identidad biológica & Una visión simbiótica de la vida: nunca hemos sido individuos & Procesos de la vida: ensayos sobre filosofía de la biología \\
\hline
\textbf{Autores y Afiliación} & Thomas Pradeu, es profesor e investigador en Filosofía de la Biología e Inmunología en la Universidad de Bordeaux / CNRS. & Scott F. Gilbert, Jan Sapp y Alfred I. Tauber, profesores e investigadores de Biología y Filosofía en Swarthmore College. & John Dupré, es profesor de Filosofía de la Ciencia en la Universidad de Exeter y director del centro EGENIS. \\
\hline
\textbf{Pertinencia} & Considero que este texto es importante porque se le planta de frente a esa visión vieja y cerrada de que los seres vivos somos organismos aislados que no dejan entrar a nadie más a su cuerpo. & Este artículo es interesante porque le hacer entender a uno que es un mito de que los animales o las personas somos entidades puras e independientes, poniendo el foco donde debe estar, en que somos más bien un equipo o un ecosistema ambulante. & Este escrito me parece relevante para pensar en cómo nombramos y entendemos las cosas, porque no podemos andar tratando a las células, plantas o animales como si fueran objetos fijos congelados en el tiempo. \\
\hline
\textbf{Introducción} & El escrito cuestiona los criterios genéticos tradicionales que fallan al intentar definir qué es un individuo, buscando demostrar que el sistema inmune no es solo un ejército defensivo, al final demuestra que en verdad se la pasa negociando y tolerando bacterias y organismos externos que son vitales para la vida. & Los autores buscan responder si de verdad se puede hablar de un "individuo" en la biología moderna cuando la ciencia demuestra que no podemos ni nacer ni digerir la comida sin microbios, su planteamiento concluye que las funciones vitales más básicas dependen por completo de las señales químicas que intercambiamos con miles de especies de bacterias. & Argumenta que la biología debe dejar de ver la vida como un conjunto de "cosas" o estructuras estáticas para empezar a entenderla como un flujo constante de procesos e información, concluye que para comprender la evolución y el funcionamiento de los seres vivos debemos enfocarnos en cómo circula la información y no en ver los organismos como cosas fijas. \\
\hline
\textbf{Importancia y aporte} & Me parece clave porque nos muestra que la identidad de un ser vivo no depende de estar encerrado en una burbuja genéticamente pura, sino de cómo se comunica y se coordina con los demás organismos que viven dentro de él. & Aporta una mirada bastante realista, porque nos recuerda que no logramos entender la anatomía ni la evolución si seguimos estudiando a los animales como entes aislados sin mirar el microbiota que los acompaña. & Es un aporte importante porque nos recuerda que el lenguaje que usamos a veces nos pone trampas, haciéndonos creer que la naturaleza está dividida en cajitas fijas cuando en realidad todo está cambiando todo el tiempo. \\
\hline
\textbf{Métodos: Datos y análisis} & Evaluó mecanismos inmunológicos analizando los patrones de reconocimiento químico que usa el cuerpo para medir hasta dónde llega lo propio (self) y lo ajeno (non-self). & Revisaron análisis de interacciones celulares y datos de metagenómica en holobiontes para medir qué tanto dependen las células del hospedador de la presencia de bacterias simbiontes. & Analizó conceptualmente la estabilidad de los sistemas vivos usando teorías de sistemas complejos para evaluar cómo se mantienen los organismos a lo largo del tiempo. \\
\hline
\textbf{Métodos: Justificación} & Prefirió irse por este camino porque los enfoques genéticos tradicionales se quedan cortos al no poder explicar por qué el cuerpo acepta y necesita a millones de microbios que no comparten su mismo ADN. & Usaron este enfoque porque sostienen que la biología tradicional cometió el error de ignorar a los microorganismos, cuando en realidad son ellos los que activan genes clave para el desarrollo del cuerpo. & Eligió esa forma de trabajo porque sostiene que la visión mecanicista tradicional de la ciencia falla al intentar explicar cómo un organismo cambia a cada segundo y aun así logra mantener su equilibrio. \\
\hline
\textbf{Resultados} & Muestra que el sistema inmune no va rechazando a ciegas todo lo que venga de afuera, sino que sabe tolerar y acoger a los simbiontes que son necesarios para sobrevivir. & Demuestra a nivel biológico y genético que un organismo no puede completar su desarrollo ni funcionar correctamente si no recibe las señales producidas por especies distintas a él. & Deja claro que las supuestas categorías fijas, como "especie" o "individuo", no son objetos estáticos reales, sino etapas o fotos fijas dentro de un proceso en constante movimiento. \\
\hline
\textbf{Tablas, leyendas y figuras} & Al ser un trabajo de reflexión filosófica e inmunológica teórica, no trae gráficas ni tablas estadísticas. & Tampoco incluye tablas ni esquemas gráficos, ya que se centra en una síntesis conceptual y revisión crítica de la literatura biológica. & Es un texto analítico y de reflexión filosófica sobre la ciencia, por lo que no lleva tablas ni esquemas gráficos. \\
\hline
\textbf{Discusión: Aporte al tema} & Nos ayuda a entender que la frontera de un individuo no es una pared fija de piel o membrana, sino una frontera inmunológica dinámica en permanente diálogo. & Nos invita a cambiar el pensamiento y reemplazar la vieja noción de "individuo genéticamente puro" por el concepto de holobionte, que entiende al ser vivo como una red de cooperación. & Conecta la manera en que funciona el lenguaje con la forma en que los científicos interpretan la información biológica y la genética. \\
\hline
\textbf{Próxima investigación} & Sería interesante investigar en qué momento exacto del desarrollo el sistema inmune decide qué bacterias extrañas va a aceptar para siempre y cuáles va a atacar. & Valdría la pena investigar cómo las alteraciones en las señales del microbiota durante la infancia podrían estar cambiando la forma en que se expresan ciertos genes en los adultos. & Se podría analizar de qué manera la enseñanza de la biología en colegios y universidades cambiaría si en vez de enseñar estructuras anatómicas fijas se enseñara la vida como redes de información en flujo. \\
\hline
\textbf{Conclusión} & Definir a un ser vivo por su aislamiento es un error, la verdadera identidad biológica se construye y se mantiene por una comunicación continua y una negociación con el entorno. & Frente a la idea clásica del individuo, la biología demuestra que los seres vivos no somos entidades independientes, sino redes interconectadas de cooperación entre múltiples especies. & Por coherencia con lo que muestra la ciencia, nos toca dejar de ver la biología como una colección de cosas estáticas; la vida no es una cosa, sino un proceso de comunicación y cambio continuo. \\
\hline
\end{longtable}

\newpage

\section*{La Lingüística de la Biología: códigos genéticos, señales celulares y comunicación en la biología}

\textbf{Autor:} E.A. Alfonso \\
\textbf{Código:} 164004902 \\
\textbf{Curso:} Elementos de la Comunicación \\
\textbf{Institución:} Universidad de los Llanos \\

La biología, durante mucho tiempo, ha sido vista como una materia experimental y descriptiva, esto debido a que en la mayoría de los colegios y universidades la muestran y la dictan como tal, como si solo fueran tubos de ensayo, reacciones químicas y listas larguísimas de nombres que toca memorizar para presentar un examen y al menos pasar, pero cuando uno se detiene a pensarlo, la vida no es solo eso, no solo es un conjunto de reacciones químicas, sino que es un sistema de comunicación, un sistema de información, un sistema de señales y códigos que se transmiten entre células, organismos y especies.

Si lo vemos desde otra perspectiva se parece bastante a cómo nos comunicamos los humanos, con palabras, gestos, señales y códigos que nos permiten transmitir y recibir información para entendernos entre nosotros. La ciencia ha encontrado la forma de expresar la información biológica en términos de códigos y señales, y ha desarrollado herramientas para entender estos códigos y cómo funcionan los sistemas biológicos, incluso de las pequeñas cosas que muchas veces no vemos a simple vista.

Por esto la idea de lenguaje se ajusta bien para comprender un poco mejor y asociar la idea de lingüística, porque un organismo no puede funcionar sin comunicarse, sin que sus partes no estén hablando todo el tiempo, organizándose y coordinándose para saber qué hacer, cuándo crecer, cuándo parar. Por esto no es raro pensar que la biología es en el fondo una ciencia de la comunicación y un constante intercambio de información, y que la lingüística es una herramienta para entender mejor cómo funciona la biología.

Cuando se piensa en biología, uno de los mayores referentes es el ADN, que es el código genético que contiene la información necesaria para construir y mantener un organismo. El ADN se compone de cuatro bases nitrogenadas (adenina, timina, citosina y guanina) que se combinan en secuencias específicas para formar genes, estos genes son como palabras en un lenguaje biológico, y su combinación determina las características de un organismo.

Entonces lo curioso es que es como un mapa o libro de la vida, aunque suena un poco raro o cliché, pero es la verdad, funciona como un escrito donde las palabras vienen siendo combinaciones de letras que, si no están conectadas entre sí o no tienen un orden, no tienen sentido.

Y es lo mismo con el ADN, una sola base nitrogenada no puede decir mucho por sí sola, pero cuando se combinan de a tres forman los codones, que vistos desde la perspectiva de la lingüística son como palabras que le dan la orden a la maquinaria celular de fabricar proteínas.

Ahí es donde la idea se ajusta y es interesante, porque si a un mensaje de texto le hace falta una letra, toda la frase puede perder sentido y terminar diciendo algo totalmente diferente a lo que se quería decir; lo mismo pasa con el ADN: si una letra de la secuencia de bases nitrogenadas se cambia, puede cambiar el mensaje que se quiere transmitir y tener consecuencias graves para el organismo.

Pero la cosa no se queda solo en tener un texto escrito y leerlo, porque en cualquier idioma una palabra puede cambiar el sentido de una frase según el tono en el que se diga o el contexto en el que se utilice, y lo mismo ocurre en biología gracias a la ciencia que se llama epigenética. (Jablonka \& Lamb, 2005).

Si el ADN es el libro de las instrucciones, la epigenética es como el entorno diciéndole a la célula cómo tiene que interpretar esa lectura, lo que demuestra que la información biológica no solo está en el código genético, sino también en cómo se interpreta y se regula esa información, la cual no está grabada en piedra, sino que puede cambiar según las condiciones del entorno.

Ahora si vamos más lejos, si salimos del núcleo de la célula y miramos el cuerpo completo, nos damos cuenta de que las células tampoco viven aisladas en su propio mundo, sino que se la pasan enviándose mensajes químicos todo el tiempo a través de receptores y moléculas de señalización, lo cual funciona igual que una red de chat o llamadas telefónicas donde hay un emisor que envía la señal, un canal por donde viaja el mensaje y un receptor que debe entenderlo para dar una respuesta, por ejemplo cuando el cuerpo necesita regular el azúcar en la sangre, se liberan señales químicas que viajan hasta los tejidos para avisarles lo que deben hacer, y si esa comunicación se interrumpe o la señal llega distorsionada, todo el sistema empieza a fallar.

En todo canal de comunicación existe además lo que conocemos como ruido, que son esas interferencias que alteran el mensaje en el camino, y eso en la biología equivale a las mutaciones o a los pequeños errores al copiar el ADN.

Aunque uno lo primero que pensaría es que un error en el mensaje es algo negativo que arruina la orden original, en los seres vivos ocurre algo curioso, porque ese ruido es precisamente el que genera la variabilidad, sin esos pequeños "fallos " en la secuencia de bases, no existirían características nuevas en las siguientes generaciones, lo que demuestra que el ruido en el canal no es solo una falla del sistema, sino la fuerza principal que le permite a las especies cambiar, adaptarse y seguir evolucionando a lo largo del tiempo.

Al final, entender la biología desde la lingüística nos ayuda a dejar de lado esa visión fría, difícil o aburrida de ver la ciencia ya que no es solo nombres difíciles que nos enseñaban en el colegio, permitiéndonos ver que la vida no es un conjunto de reacciones químicas ocurriendo por suerte o al azar, sino una red en permanente conversación donde cada molécula, cada gen y cada señal aportan su parte para que el mensaje de la vida siga fluyendo.

\end{document}\documentclass[12pt]{article}
\usepackage[utf8]{inputenc}
\usepackage[T1]{fontenc}
\usepackage[spanish]{babel}
\usepackage[letterpaper, margin=2.5cm]{geometry}
\usepackage{longtable}
\usepackage{array}

\begin{document}

\section*{Tabla 1: Artículos de Actividad de Clase}

\begin{longtable}{|>{\raggedright\arraybackslash}p{3cm}|>{\raggedright\arraybackslash}p{4.2cm}|>{\raggedright\arraybackslash}p{4.2cm}|>{\raggedright\arraybackslash}p{4.2cm}|}
\hline
\textbf{Criterio} & \textbf{Artículo 1} & \textbf{Artículo 2} & \textbf{Artículo 3} \\
\hline
\endfirsthead

\hline
\textbf{Criterio} & \textbf{Artículo 1} & \textbf{Artículo 2} & \textbf{Artículo 3} \\
\hline
\endhead

\textbf{Título} & Cambio climático y responsabilidad individual: Una respuesta a Johnson & Obligaciones éticas en una tragedia de los bienes comunes & Clima, acción colectiva y obligaciones morales individuales \\
\hline
\textbf{Autores y Afiliación} & Marion Hourdequin, es profesora de Filosofía en Colorado College, con trayectoria e investigación en ética ambiental. & Baylor L. Johnson, profesor de Filosofía en St. Lawrence University. & Marion Hourdequin, profesora de Filosofía en Colorado College. \\
\hline
\textbf{Pertinencia} & Considero que este texto es muy importante porque se le planta de frente a esa visión pesimista de que las acciones que uno hace a diario en la casa no sirven de nada frente a un problema tan gigante como el cambio climático. & Este artículo es interesante porque evita que toda la culpa del problema caiga encima del consumidor, poniendo el foco donde debe estar, en la responsabilidad de las instituciones y del Estado. & Este escrito me parece relevante para pensar en la coherencia de nuestras acciones, porque no podemos andar pidiéndole cambios drásticos al gobierno o a las empresas si nosotros en la casa no asumimos ni un poco de responsabilidad. \\
\hline
\textbf{Introducción} & El escrito le responde directamente a Johnson para defender que lo que hacemos individualmente sí tiene peso, buscando replantear la famosa "tragedia de los bienes comunes", al final demuestra que cuando uno reduce su huella personal, no solo está haciendo un cambio chiquito, sino les manda un mensaje ético y claro a los demás y los motiva a actuar también. & Johnson busca responder si de verdad estamos obligados a hacer sacrificios personales cuando el problema es colectivo y gigante, su planteamiento concluye que nuestra verdadera responsabilidad no es sufrir solos, sino presionar y apoyar acuerdos que apliquen para todo el mundo por igual. & Argumenta que sí tenemos el deber moral de bajar nuestras emisiones, no por salvar el planeta solos, sino por un tema de integridad personal y de respeto con las relaciones que construimos con los demás, concluye que, para tener la moral de exigir normas públicas primero debemos ser coherentes en la vida privada de cada uno. \\
\hline
\textbf{Importancia y aporte} & Me parece clave porque nos muestra que el compromiso individual y las políticas públicas tienen que ir cogidos de la mano, no es que uno reemplace al otro ni que sean cosas opuestas. & Aporta una mirada bastante realista, porque nos recuerda que no logramos mucho cambiando hábitos a puerta cerrada si al mismo tiempo no estamos exigiéndole a leyes e instituciones firmes. & Es un aporte importante porque nos recuerda que la protesta y la exigencia política cogen más fuerza y credibilidad cuando están respaldadas con el propio ejemplo. \\
\hline
\textbf{Métodos: Datos y análisis} & Analizó conductas de todos los días, como usar la bicicleta o ponerse a reciclar, y cómo eso impacta en el círculo cercano. & Usó escenarios teóricos para analizar cómo se comporta un grupo de gente cuando comparte un recurso común pero no hay reglas claras de por medio. & Analizó conceptualmente cómo nos conectamos con la comunidad y de qué manera lo que hacemos a diario termina tocando esos vínculos sociales. \\
\hline
\textbf{Métodos: Justificación} & Prefirió irse por este camino porque la teoría de juegos tradicional siempre asume que la gente actúa pura y exclusivamente por egoísmo, mientras que los enfoques tipo Ostrom reconocen que a las personas también nos mueven los valores y las normas sociales. & Uso este modelo porque argumenta que, si no existe una ley general, el esfuerzo que uno hace con tanto cariño termina siendo aplastado por el consumo desmedido de los demás. & Eligió esa forma de trabajo porque sostiene que los humanos no somos seres aislados, sino seres interconectados que nos influenciamos mutuamente todo el tiempo. \\
\hline
\textbf{Resultados} & Muestra que lo que hacemos en el día a día no es del todo invisible, sino que se vuelve una especie de mensaje o señal que termina influyendo en la forma de actuar de la gente que nos rodea. & Demuestra a nivel teórico que hacer sacrificios aislados sin obligar a los demás a cumplir las mismas reglas no soluciona el problema de fondo. & Deja claro que cuando uno dice una cosa, pero hace otra en la vida diaria, pierde totalmente la autoridad moral a la hora de exigirle cambios reales a las instituciones. \\
\hline
\textbf{Tablas, leyendas y figuras} & Al ser un trabajo de reflexión filosófica pura y dura, no trae gráficas ni tablas estadísticas. & Tampoco incluye tablas ni esquemas gráficos, ya que su desarrollo es de análisis teórico. & Es un texto analítico y de reflexión moral, por lo que no lleva tablas ni esquemas gráficos. \\
\hline
\textbf{Discusión: Aporte al tema} & Nos ayuda a entender que las leyes y las grandes políticas no se sostienen solas en el aire, sino que necesitan de la confianza y los valores de la propia comunidad para que realmente funcionen. & Nos invita a cambiar el pensamiento y pasar del simple consumo responsable, al activismo político y ciudadano. & Conecta lo que hacemos a puerta cerrada en nuestra intimidad con nuestro papel como ciudadanos activos en la sociedad. \\
\hline
\textbf{Próxima investigación} & Sería interesante investigar cuánta gente en un barrio o comunidad necesita adoptar hábitos sostenibles antes de que las autoridades locales hagan el deber de aprobar una normativa ambiental obligatoria. & Valdría la pena investigar qué clase de incentivos o argumentos harían que los ciudadanos voten a favor de leyes ambientales mucho más estrictas, aunque impliquen cambios difíciles en su vida diaria. & Se podría analizar cómo la percepción de coherencia en los líderes ambientales influye en la disposición de la gente a seguir sus propuestas o apoyarlos en las urnas. \\
\hline
\textbf{Conclusión} & Actuar bien por decisión propia no va a solucionar el problema global por sí solo, pero sí ayuda a construir ese piso de confianza necesario para que las leyes e instituciones puedan dar resultados. & Frente a un desastre global, nuestra principal obligación ética no es andar haciendo sacrificios solos, sino impulsar y exigir reglas colectivas que nos obliguen a todos a cumplir por igual dejando la desigualdad de lado. & Por integridad nos toca reducir la huella personal, no porque con eso arreglemos el clima por arte de magia, sino porque le da un respaldo ético indestructible a cualquier cambio político que salgamos a exigir. \\
\hline
\end{longtable}

\newpage

\section*{Tabla 2: Matriz de Lectura - Una biología sin individuos ontológicos}

\begin{longtable}{|>{\raggedright\arraybackslash}p{3cm}|>{\raggedright\arraybackslash}p{4.2cm}|>{\raggedright\arraybackslash}p{4.2cm}|>{\raggedright\arraybackslash}p{4.2cm}|}
\hline
\textbf{Criterio} & \textbf{Artículo 1} & \textbf{Artículo 2} & \textbf{Artículo 3} \\
\hline
\endfirsthead

\hline
\textbf{Criterio} & \textbf{Artículo 1} & \textbf{Artículo 2} & \textbf{Artículo 3} \\
\hline
\endhead

\textbf{Título} & Los límites del yo: inmunología e identidad biológica & Una visión simbiótica de la vida: nunca hemos sido individuos & Procesos de la vida: ensayos sobre filosofía de la biología \\
\hline
\textbf{Autores y Afiliación} & Thomas Pradeu, es profesor e investigador en Filosofía de la Biología e Inmunología en la Universidad de Bordeaux / CNRS. & Scott F. Gilbert, Jan Sapp y Alfred I. Tauber, profesores e investigadores de Biología y Filosofía en Swarthmore College. & John Dupré, es profesor de Filosofía de la Ciencia en la Universidad de Exeter y director del centro EGENIS. \\
\hline
\textbf{Pertinencia} & Considero que este texto es importante porque se le planta de frente a esa visión vieja y cerrada de que los seres vivos somos organismos aislados que no dejan entrar a nadie más a su cuerpo. & Este artículo es interesante porque le hacer entender a uno que es un mito de que los animales o las personas somos entidades puras e independientes, poniendo el foco donde debe estar, en que somos más bien un equipo o un ecosistema ambulante. & Este escrito me parece relevante para pensar en cómo nombramos y entendemos las cosas, porque no podemos andar tratando a las células, plantas o animales como si fueran objetos fijos congelados en el tiempo. \\
\hline
\textbf{Introducción} & El escrito cuestiona los criterios genéticos tradicionales que fallan al intentar definir qué es un individuo, buscando demostrar que el sistema inmune no es solo un ejército defensivo, al final demuestra que en verdad se la pasa negociando y tolerando bacterias y organismos externos que son vitales para la vida. & Los autores buscan responder si de verdad se puede hablar de un "individuo" en la biología moderna cuando la ciencia demuestra que no podemos ni nacer ni digerir la comida sin microbios, su planteamiento concluye que las funciones vitales más básicas dependen por completo de las señales químicas que intercambiamos con miles de especies de bacterias. & Argumenta que la biología debe dejar de ver la vida como un conjunto de "cosas" o estructuras estáticas para empezar a entenderla como un flujo constante de procesos e información, concluye que para comprender la evolución y el funcionamiento de los seres vivos debemos enfocarnos en cómo circula la información y no en ver los organismos como cosas fijas. \\
\hline
\textbf{Importancia y aporte} & Me parece clave porque nos muestra que la identidad de un ser vivo no depende de estar encerrado en una burbuja genéticamente pura, sino de cómo se comunica y se coordina con los demás organismos que viven dentro de él. & Aporta una mirada bastante realista, porque nos recuerda que no logramos entender la anatomía ni la evolución si seguimos estudiando a los animales como entes aislados sin mirar el microbiota que los acompaña. & Es un aporte importante porque nos recuerda que el lenguaje que usamos a veces nos pone trampas, haciéndonos creer que la naturaleza está dividida en cajitas fijas cuando en realidad todo está cambiando todo el tiempo. \\
\hline
\textbf{Métodos: Datos y análisis} & Evaluó mecanismos inmunológicos analizando los patrones de reconocimiento químico que usa el cuerpo para medir hasta dónde llega lo propio (self) y lo ajeno (non-self). & Revisaron análisis de interacciones celulares y datos de metagenómica en holobiontes para medir qué tanto dependen las células del hospedador de la presencia de bacterias simbiontes. & Analizó conceptualmente la estabilidad de los sistemas vivos usando teorías de sistemas complejos para evaluar cómo se mantienen los organismos a lo largo del tiempo. \\
\hline
\textbf{Métodos: Justificación} & Prefirió irse por este camino porque los enfoques genéticos tradicionales se quedan cortos al no poder explicar por qué el cuerpo acepta y necesita a millones de microbios que no comparten su mismo ADN. & Usaron este enfoque porque sostienen que la biología tradicional cometió el error de ignorar a los microorganismos, cuando en realidad son ellos los que activan genes clave para el desarrollo del cuerpo. & Eligió esa forma de trabajo porque sostiene que la visión mecanicista tradicional de la ciencia falla al intentar explicar cómo un organismo cambia a cada segundo y aun así logra mantener su equilibrio. \\
\hline
\textbf{Resultados} & Muestra que el sistema inmune no va rechazando a ciegas todo lo que venga de afuera, sino que sabe tolerar y acoger a los simbiontes que son necesarios para sobrevivir. & Demuestra a nivel biológico y genético que un organismo no puede completar su desarrollo ni funcionar correctamente si no recibe las señales producidas por especies distintas a él. & Deja claro que las supuestas categorías fijas, como "especie" o "individuo", no son objetos estáticos reales, sino etapas o fotos fijas dentro de un proceso en constante movimiento. \\
\hline
\textbf{Tablas, leyendas y figuras} & Al ser un trabajo de reflexión filosófica e inmunológica teórica, no trae gráficas ni tablas estadísticas. & Tampoco incluye tablas ni esquemas gráficos, ya que se centra en una síntesis conceptual y revisión crítica de la literatura biológica. & Es un texto analítico y de reflexión filosófica sobre la ciencia, por lo que no lleva tablas ni esquemas gráficos. \\
\hline
\textbf{Discusión: Aporte al tema} & Nos ayuda a entender que la frontera de un individuo no es una pared fija de piel o membrana, sino una frontera inmunológica dinámica en permanente diálogo. & Nos invita a cambiar el pensamiento y reemplazar la vieja noción de "individuo genéticamente puro" por el concepto de holobionte, que entiende al ser vivo como una red de cooperación. & Conecta la manera en que funciona el lenguaje con la forma en que los científicos interpretan la información biológica y la genética. \\
\hline
\textbf{Próxima investigación} & Sería interesante investigar en qué momento exacto del desarrollo el sistema inmune decide qué bacterias extrañas va a aceptar para siempre y cuáles va a atacar. & Valdría la pena investigar cómo las alteraciones en las señales del microbiota durante la infancia podrían estar cambiando la forma en que se expresan ciertos genes en los adultos. & Se podría analizar de qué manera la enseñanza de la biología en colegios y universidades cambiaría si en vez de enseñar estructuras anatómicas fijas se enseñara la vida como redes de información en flujo. \\
\hline
\textbf{Conclusión} & Definir a un ser vivo por su aislamiento es un error, la verdadera identidad biológica se construye y se mantiene por una comunicación continua y una negociación con el entorno. & Frente a la idea clásica del individuo, la biología demuestra que los seres vivos no somos entidades independientes, sino redes interconectadas de cooperación entre múltiples especies. & Por coherencia con lo que muestra la ciencia, nos toca dejar de ver la biología como una colección de cosas estáticas; la vida no es una cosa, sino un proceso de comunicación y cambio continuo. \\
\hline
\end{longtable}

\newpage

\section*{La Lingüística de la Biología: códigos genéticos, señales celulares y comunicación en la biología}

\textbf{Autor:} E.A. Alfonso \\
\textbf{Código:} 164004902 \\
\textbf{Curso:} Elementos de la Comunicación \\
\textbf{Institución:} Universidad de los Llanos \\

La biología, durante mucho tiempo, ha sido vista como una materia experimental y descriptiva, esto debido a que en la mayoría de los colegios y universidades la muestran y la dictan como tal, como si solo fueran tubos de ensayo, reacciones químicas y listas larguísimas de nombres que toca memorizar para presentar un examen y al menos pasar, pero cuando uno se detiene a pensarlo, la vida no es solo eso, no solo es un conjunto de reacciones químicas, sino que es un sistema de comunicación, un sistema de información, un sistema de señales y códigos que se transmiten entre células, organismos y especies.

Si lo vemos desde otra perspectiva se parece bastante a cómo nos comunicamos los humanos, con palabras, gestos, señales y códigos que nos permiten transmitir y recibir información para entendernos entre nosotros. La ciencia ha encontrado la forma de expresar la información biológica en términos de códigos y señales, y ha desarrollado herramientas para entender estos códigos y cómo funcionan los sistemas biológicos, incluso de las pequeñas cosas que muchas veces no vemos a simple vista.

Por esto la idea de lenguaje se ajusta bien para comprender un poco mejor y asociar la idea de lingüística, porque un organismo no puede funcionar sin comunicarse, sin que sus partes no estén hablando todo el tiempo, organizándose y coordinándose para saber qué hacer, cuándo crecer, cuándo parar. Por esto no es raro pensar que la biología es en el fondo una ciencia de la comunicación y un constante intercambio de información, y que la lingüística es una herramienta para entender mejor cómo funciona la biología.

Cuando se piensa en biología, uno de los mayores referentes es el ADN, que es el código genético que contiene la información necesaria para construir y mantener un organismo. El ADN se compone de cuatro bases nitrogenadas (adenina, timina, citosina y guanina) que se combinan en secuencias específicas para formar genes, estos genes son como palabras en un lenguaje biológico, y su combinación determina las características de un organismo.

Entonces lo curioso es que es como un mapa o libro de la vida, aunque suena un poco raro o cliché, pero es la verdad, funciona como un escrito donde las palabras vienen siendo combinaciones de letras que, si no están conectadas entre sí o no tienen un orden, no tienen sentido.

Y es lo mismo con el ADN, una sola base nitrogenada no puede decir mucho por sí sola, pero cuando se combinan de a tres forman los codones, que vistos desde la perspectiva de la lingüística son como palabras que le dan la orden a la maquinaria celular de fabricar proteínas.

Ahí es donde la idea se ajusta y es interesante, porque si a un mensaje de texto le hace falta una letra, toda la frase puede perder sentido y terminar diciendo algo totalmente diferente a lo que se quería decir; lo mismo pasa con el ADN: si una letra de la secuencia de bases nitrogenadas se cambia, puede cambiar el mensaje que se quiere transmitir y tener consecuencias graves para el organismo.

Pero la cosa no se queda solo en tener un texto escrito y leerlo, porque en cualquier idioma una palabra puede cambiar el sentido de una frase según el tono en el que se diga o el contexto en el que se utilice, y lo mismo ocurre en biología gracias a la ciencia que se llama epigenética. (Jablonka \& Lamb, 2005).

Si el ADN es el libro de las instrucciones, la epigenética es como el entorno diciéndole a la célula cómo tiene que interpretar esa lectura, lo que demuestra que la información biológica no solo está en el código genético, sino también en cómo se interpreta y se regula esa información, la cual no está grabada en piedra, sino que puede cambiar según las condiciones del entorno.

Ahora si vamos más lejos, si salimos del núcleo de la célula y miramos el cuerpo completo, nos damos cuenta de que las células tampoco viven aisladas en su propio mundo, sino que se la pasan enviándose mensajes químicos todo el tiempo a través de receptores y moléculas de señalización, lo cual funciona igual que una red de chat o llamadas telefónicas donde hay un emisor que envía la señal, un canal por donde viaja el mensaje y un receptor que debe entenderlo para dar una respuesta, por ejemplo cuando el cuerpo necesita regular el azúcar en la sangre, se liberan señales químicas que viajan hasta los tejidos para avisarles lo que deben hacer, y si esa comunicación se interrumpe o la señal llega distorsionada, todo el sistema empieza a fallar.

En todo canal de comunicación existe además lo que conocemos como ruido, que son esas interferencias que alteran el mensaje en el camino, y eso en la biología equivale a las mutaciones o a los pequeños errores al copiar el ADN.

Aunque uno lo primero que pensaría es que un error en el mensaje es algo negativo que arruina la orden original, en los seres vivos ocurre algo curioso, porque ese ruido es precisamente el que genera la variabilidad, sin esos pequeños "fallos " en la secuencia de bases, no existirían características nuevas en las siguientes generaciones, lo que demuestra que el ruido en el canal no es solo una falla del sistema, sino la fuerza principal que le permite a las especies cambiar, adaptarse y seguir evolucionando a lo largo del tiempo.

Al final, entender la biología desde la lingüística nos ayuda a dejar de lado esa visión fría, difícil o aburrida de ver la ciencia ya que no es solo nombres difíciles que nos enseñaban en el colegio, permitiéndonos ver que la vida no es un conjunto de reacciones químicas ocurriendo por suerte o al azar, sino una red en permanente conversación donde cada molécula, cada gen y cada señal aportan su parte para que el mensaje de la vida siga fluyendo.

\end{document}

